\documentclass[11pt]{article}

\usepackage[margin=1in]{geometry}
\usepackage{amsmath,amssymb,booktabs,tabularx}

\title{Affine domain constraints}
\author{}
\date{}

\begin{document}
\maketitle

\noindent\textbf{Status.} Draft\\
\textbf{Plan item.} Affine domain constraints\\
\textbf{Depends on.} \texttt{02\_first\_moment\_guard.tex},
\texttt{04\_affine\_transform\_ansatz.tex},
\texttt{05\_riccati\_ode\_system\_for\_c\_d\_v.tex},
\texttt{15\_fourier\_business\_time\_clock\_interval\_factor\_grid.tex}\\
\textbf{Feeds.} Affine schedule transform, terminal pricing measure, admissibility checks, benchmark governance

\section{Purpose}

State the mathematical domain on which the affine transform is finite and the
ODE system is well-defined.
These constraints separate admissible transform queries from parameter or
argument choices that would make the conditional moment explode or leave the
model's state domain.

\section{Setup}

Use the affine schedule transform from
\texttt{04\_affine\_transform\_ansatz.tex} and the ODE system from
\texttt{05\_riccati\_ode\_system\_for\_c\_d\_v.tex}.
The domain is checked on the whole backward integration interval
\(\tau\in[0,\Delta]\), not only at the terminal or initial endpoint.

\section{Notation}

\begin{tabularx}{\textwidth}{lXX}
\toprule
Symbol & Meaning & Units or domain\\
\midrule
\(\mathcal D_{\mathrm{aff}}\) & Admissible affine transform domain & Set of parameters, state, and arguments\\
\(g_c(\tau)\) & Gamma-normal news domain metric for component \(c\) & Complex\\
\(\theta_U\) & Unit-mean gamma severity scale & Strictly positive\\
\(k_U\) & Gamma severity shape & Strictly positive\\
\(\rho(\mathbf A_\alpha/\beta)\) & Hawkes branching spectral radius & Strictly less than one\\
\(\bar z\) & Derived long-variance mean-reversion level & Strictly positive\\
\(p_B^\pm,p_E^\pm\) & Baseline-source and endogenous-source move probabilities & Unit interval\\
\(C,\mathbf D,V\) & Affine ODE solution & Finite on \([0,\Delta]\)\\
\bottomrule
\end{tabularx}

\section{Assumptions used}

\begin{tabularx}{\textwidth}{XXX}
\toprule
Assumption & Where it is used & Why it is needed\\
\midrule
Hawkes parameters satisfy non-negativity and branching stability.
& Defining valid intensities and stationary rates for the long-variance normalization.
& The model needs non-negative intensities and finite first moments.\\
The long-variance first-moment guard holds.
& Ensuring \(\bar z>0\).
& The CIR mean-reversion level must be positive under unit-mean normalization.\\
News mark transforms are evaluated only inside the gamma MGF domain.
& Defining \(M_c\) and \(\ell_c\).
& The gamma-normal transform is finite only before the gamma pole.\\
The affine ODE solution is finite throughout integration.
& Evaluating \(F=\exp(C+\mathbf D\cdot\mathbf y_0+VZ_0)\).
& A finite-time Riccati or Hamiltonian blowup invalidates the transform.\\
Boundary states are inside the model state domain.
& Evaluating the final affine exponent.
& The transform is conditional on admissible Hawkes and long-variance states.\\
\bottomrule
\end{tabularx}

\section{Result}

\noindent\textbf{Proposition.} A schedule-interval affine transform query is
admissible if all of the following hold.

\begin{enumerate}
    \item The Hawkes kernel satisfies
    \[
        \bar\lambda>0,\qquad
        \beta>0,\qquad
        \alpha_{ij}\ge 0,\qquad
        \rho_{ij}\ge 0,
        \qquad
        \rho(\mathbf A_\alpha/\beta)<1.
    \]
    \item The long-variance parameters satisfy
    \[
        \kappa_z>0,\qquad
        \sigma_z\ge0,\qquad
        \chi_+\ge0,\qquad
        \chi_-\ge0,
    \]
    and the first-moment guard
    \[
        \kappa_z
        >
        \chi_+\lambda_{\mathrm{stat}}^+
        +\chi_-\lambda_{\mathrm{stat}}^-,
    \]
    so that
    \[
        \bar z
        =
        1-
            \frac{
            \chi_+\lambda_{\mathrm{stat}}^+
            +\chi_-\lambda_{\mathrm{stat}}^-
        }{\kappa_z}
        >0.
    \]
    \item The boundary state satisfies
    \[
        \mathbf y_0\ge0 \quad\hbox{componentwise},
        \qquad
        Z_0>0.
    \]
    \item The interval and deterministic interval-grid intensities satisfy
    \[
        \Delta>0,\qquad s_k\ge0,\qquad
        \nu_k^{\mathrm{known}}\ge0,\qquad
        \nu_k^{\mathrm{surprise}}\ge0.
    \]
    \item The terminal loadings and transform arguments are finite complex
    values, and the ODE solution \((C(\tau),\mathbf D(\tau),V(\tau))\) remains
    finite for every \(\tau\in[0,\Delta]\).
    \item For every active news component \(c\) and every
    \(\tau\in[0,\Delta]\),
    \[
        \operatorname{Re} g_c(\tau)<\frac{1}{\theta_U},
    \]
    where
    \[
        g_c(\tau)
        =
        (\mathbf D(\tau)\cdot\mathbf r)a_c
        +qm_c+\frac{1}{2}q^2\sigma_c^2.
    \]
\end{enumerate}

Under these conditions, the affine transform is finite and equals
\[
    \exp\!\left(C(\Delta)+\mathbf D(\Delta)\cdot\mathbf y_0+V(\Delta)Z_0\right).
\]

\section{Derivation}

The Hawkes conditions guarantee non-negative event intensities:
\[
    \lambda_t^+
    =
    s_k\bar\lambda+Y_t^{++}+Y_t^{+-},
    \qquad
    \lambda_t^-
    =
    s_k\bar\lambda+Y_t^{-+}+Y_t^{--}.
\]
Here \(s_k\) is the non-negative interval factor materialized from the
Fourier business-time clock for the schedule interval.
The spectral-radius condition is the finite-branching condition for the
linear Hawkes reaction system.
It is also needed when stationary side rates \(\lambda_{\mathrm{stat}}^\pm\) are used to
derive the long-variance
mean-reversion level.

The long-variance normalization sets the stationary mean target to one.
Taking first moments in
\[
    dZ_t
    =
    \kappa_z(\bar z-Z_t)\,dt
    +\sigma_z\sqrt{Z_t}\,dW_t^Z
    +\chi_+\,dN_t^+
    +\chi_-\,dN_t^-
\]
gives
\[
    \bar z
    =
    1-
    \frac{
        \chi_+\lambda_{\mathrm{stat}}^+
        +\chi_-\lambda_{\mathrm{stat}}^-
    }{\kappa_z}.
\]
The first-moment guard is exactly the condition \(\bar z>0\).

For news marks, condition on a component-\(c\) arrival.
The activity and return mark multiplier is
\[
    \mathbb E\!\left[
        \exp\!\left(g_c(\tau)U\right)
    \right],
\]
after the conditional normal return shock has been integrated out.
Since \(U\sim\Gamma(k_U,\theta_U)\), this expectation is
\[
    M_c(q,\mathbf D(\tau)\cdot\mathbf r)
    =
    \left(1-\theta_U g_c(\tau)\right)^{-k_U},
\]
which is finite on the moment half-plane
\[
    \operatorname{Re}g_c(\tau)<\frac{1}{\theta_U}.
\]
This condition also keeps \(1-\theta_Ug_c(\tau)\) in the right half-plane, so
the principal power and principal logarithm used for \(M_c\) and \(\ell_c\)
are single-valued along the integration path.

The remaining domain requirement is finite ODE flow.
The \(V\) equation is a Riccati equation:
\[
    V'
    =
    -\kappa_zV
    +\frac{1}{2}\sigma_z^2 V^2
    -s_k\bar\lambda(p_B^+u_+^B+p_B^-u_-^B).
\]
If this equation develops a finite-time pole before \(\Delta\), then the
exponent is not finite.
The Hawkes loading equations also contain exponential Hamiltonians,
so non-finite \(D\) or \(V\) invalidates the event multipliers.
Therefore \(C\), \(\mathbf D\), and \(V\) must stay finite on the full
integration interval.

\section{Domain and validity}

The proposition states the full interval domain.
Components with zero news intensity do not constrain the active news moment
domain, because they do not enter the interval transform.
The square-root long-variance process is used on its non-negative affine
diffusion domain with positive boundary state and positive derived
mean-reversion level.
No additional Feller condition is needed for affine closure.

The condition \(\operatorname{Re}g_c(\tau)<1/\theta_U\) is strict.
Equality is not admissible because it is the gamma pole.
Values beyond the boundary do not define finite moments.

\section{Consequences}

\begin{itemize}
    \item Pricing and paper plots must report failures as domain failures when the requested moment lies outside the affine transform domain.
    \item The news-domain guard is pathwise in ODE time, not a check only on the final exponent.
    \item The first-moment guard is a parameter-domain condition, while Riccati blowup is a transform-argument condition.
    \item Inactive news components impose no moment-domain restriction on an interval.
\end{itemize}

\section{Checks}

\begin{tabularx}{\textwidth}{XXX}
\toprule
Check & Expected result & Evidence location\\
\midrule
News-domain metric
& The guard metric is \(g_c=(\mathbf D\cdot\mathbf r)a_c+qm_c+\frac12q^2\sigma_c^2\), selected over active news components.
& \texttt{gamma\_normal\_news\_domain\_metric} in \texttt{crates/janiform-affine-oracle/src/affine\_oracle/news\_transform.rs}\\
News-domain upper bound
& Guarded integration terminates when the real metric reaches \(1/\theta_U\).
& Test \texttt{solver\_domain\_guard\_violation\_terminates\_with\_deterministic\_metadata}\\
Schedule checkpoint guard
& Schedule integration checks finite \(C\), finite \(\mathbf D\), finite \(V\), and the news-domain metric at every checkpoint.
& \texttt{evaluate\_schedule\_guard\_checkpoint} in \texttt{crates/janiform-affine-oracle/src/affine\_oracle/schedule/core.rs}\\
Long-variance first-moment domain
& The derived mean-reversion level must be positive.
& Test \texttt{stationary\_mean\_reversion\_level\_rejects\_exhausted\_event\_feedback\_budget}\\
Input state domain
& Hawkes state components are non-negative finite values and the long-variance boundary state is positive.
& \texttt{validate\_initial\_state} in \texttt{crates/janiform-affine-oracle/src/affine\_oracle/contracts.rs} and \texttt{LongVarianceState::validate}\\
\bottomrule
\end{tabularx}

\end{document}
